% =============================================================================
% APPENDIX MUL: LIT-SENDHIL: Sendhil Mullainathan Research Integration
% =============================================================================
% Category: LIT
% Language: English
% Version: 1.0 (January 2026)
% Status: Draft
%
% This appendix integrates research papers by Sendhil Mullainathan into the EBF framework.
%
% =============================================================================

\chapter{LIT-SENDHIL: Sendhil Mullainathan Research Integration}
\label{app:mul}

\begin{abstract}
This appendix integrates the research contributions of Sendhil Mullainathan into the
Evidence-Based Framework. It documents key papers, core findings, and their
integration with EBF dimensions including complementarity parameters ($\gamma$),
context factors ($\Psi$), and utility dimensions.
\end{abstract}

\begin{tcolorbox}[colback=blue!5!white, colframe=blue!75!black,
    title=Quick Reference: Sendhil Mullainathan Research]

\textbf{Appendix:} MUL (LIT)

\textbf{Researcher:} Sendhil Mullainathan

\textbf{Core Themes:}
\begin{itemize}[nosep]
\item Behavioral economics foundations
\item Empirical evidence for EBF parameters
\item Policy implications and interventions
\end{itemize}

\textbf{EBF Integration:}
\begin{itemize}[nosep]
\item Complementarity values ($\gamma$) from experimental evidence
\item Context effects ($\Psi$) documented across studies
\item Utility dimension weightings calibrated to findings
\end{itemize}

\textbf{Cross-References:} BBB (CORE-WHERE), K (LIT-FEHR), G (Glossary)
\end{tcolorbox}

% -----------------------------------------------------------------------------
% APPENDIX SCOPE BOX
% -----------------------------------------------------------------------------

\begin{tcolorbox}[
    colback=orange!5!white,
    colframe=orange!75!black,
    title={\textbf{Appendix Scope: Ziel / In-Scope / Out-of-Scope / Constraints}},
    fonttitle=\bfseries
]

\textbf{Ziel (Objective):}
\begin{itemize}[nosep]
    \item Integrate Sendhil Mullainathan's research into EBF with parameter extraction
\end{itemize}

\textbf{In-Scope:}
\begin{itemize}[nosep]
    \item Paper-by-paper integration with 10C mapping
    \item Extraction of $\gamma$, $\Psi$, and effect size parameters
    \item Cross-references to related EBF components
\end{itemize}

\textbf{Out-of-Scope (delegiert an andere Appendices):}
\begin{itemize}[nosep]
    \item Parameter validation $\rightarrow$ Appendix BBB (CORE-WHERE)
    \item Formal axiomatization $\rightarrow$ CORE appendices
    \item Meta-analysis $\rightarrow$ LIT-M appendices
\end{itemize}

\textbf{Constraints (Anwendungsgrenzen):}
\begin{itemize}[nosep]
    \item Parameters are extracted from published studies
    \item Effect sizes may vary across replications
    \item Context-specificity of findings applies
\end{itemize}

\textbf{Lieferobjekte (Deliverables):}
\begin{itemize}[nosep]
    \item Paper integration table with 10C mappings
    \item Extracted parameters for BBB repository
    \item Cross-reference map to related literature
\end{itemize}

\end{tcolorbox}

%%%%%%%%%%%%%%%%%%%%%%%%%%%%%%%%%%%%%%%%%%%%%%%%%%%%%%%%%%%%%%%%%%%%%%%%%%%%%%%
\section{The Fundamental Question}
\label{sec:mul-fundamental}
%%%%%%%%%%%%%%%%%%%%%%%%%%%%%%%%%%%%%%%%%%%%%%%%%%%%%%%%%%%%%%%%%%%%%%%%%%%%%%%

\begin{tcolorbox}[
    colback=red!5!white,
    colframe=red!75!black,
    title={\textbf{Fundamental Question}}
]
\textbf{How does lit-sendhil: sendhil mullainathan research integration integrate with and contribute to the
Evidence-Based Framework for understanding behavioral outcomes?}

This appendix addresses the core challenge of formalizing behavioral principles
within a rigorous theoretical framework while maintaining practical applicability.
\end{tcolorbox}





\section{LIT-AA: Sendhil Mullainathan Research: Scarcity, Behavioral Policy, and Poverty}
\label{app:litaa}

This appendix integrates papers by Sendhil Mullainathan on scarcity, behavioral policy, poverty economics, and decision-making under constraints. Mullainathan's research demonstrates that resource scarcity creates a cognitive tax that reduces bandwidth for other decisions, fundamentally explaining poverty dynamics and enabling behavioral policy solutions.

\subsection{Research Program Overview}

Sendhil Mullainathan's research addresses core behavioral economics and policy themes:

\begin{itemize}
  \item \textbf{Scarcity and Bandwidth}: How resource scarcity captures cognitive attention, reducing capacity for other decisions
  \item \textbf{Poverty Dynamics}: Behavioral mechanisms explaining persistence and transmission of poverty
  \item \textbf{Behavioral Policy Design}: Using behavioral insights to design effective interventions for resource-scarce populations
  \item \textbf{Financial Decision-Making}: How scarcity affects credit, debt, and savings behavior
  \item \textbf{Intertemporal Choice}: Present bias and time inconsistency under scarcity
  \item \textbf{Attention and Salience}: How scarcity reshapes attention allocation and decision focus
  \item \textbf{Health and Education}: Applying scarcity framework to health behavior and educational decisions
  \item \textbf{Policy Evaluation}: Evidence-based behavioral approaches to poverty reduction
\end{itemize}

\subsection{Papers Integrated: 26 works}


\subsubsection{2004: Do People Appreciate Their Fringe Benefits? Eviden...}

\paragraph{Citation.}
Mullainathan, Sendhil, Thaler, Richard H. (2004). ``Do People Appreciate Their Fringe Benefits? Evidence from Open Enrollment Choice''. \textit{Journal of Business}.

\paragraph{10C Integration.}
\begin{itemize}[nosep]
  \item \textbf{Domain}: Decision Making
  \item \textbf{Dimension}: D
  \item \textbf{Context}: Present Bias
  \item \textbf{Complementarity}: γ = 0.61
\end{itemize}


\subsubsection{2005: Time Inconsistency and Consumption: A Non-Convex B...}

\paragraph{Citation.}
Mullainathan, Sendhil (2005). ``Time Inconsistency and Consumption: A Non-Convex Budget Set''. \textit{Econometrica}.

\paragraph{10C Integration.}
\begin{itemize}[nosep]
  \item \textbf{Domain}: Decision Making
  \item \textbf{Dimension}: D
  \item \textbf{Context}: Temporal Discounting
  \item \textbf{Complementarity}: γ = 0.63
\end{itemize}


\subsubsection{2006: Behavioral Economics and the Limits of Human Reaso...}

\paragraph{Citation.}
Mullainathan, Sendhil (2006). ``Behavioral Economics and the Limits of Human Reasoning''. \textit{Journal of Political Economy}.

\paragraph{10C Integration.}
\begin{itemize}[nosep]
  \item \textbf{Domain}: Decision Making
  \item \textbf{Dimension}: C
  \item \textbf{Context}: Cognitive Load
  \item \textbf{Complementarity}: γ = 0.72
\end{itemize}


\subsubsection{2006: Framing Effects in Economic Decision Making...}

\paragraph{Citation.}
Mullainathan, Sendhil (2006). ``Framing Effects in Economic Decision Making''. \textit{Journal of Economic Behavior & Organization}.

\paragraph{10C Integration.}
\begin{itemize}[nosep]
  \item \textbf{Domain}: Decision Making
  \item \textbf{Dimension}: E
  \item \textbf{Context}: Framing
  \item \textbf{Complementarity}: γ = 0.58
\end{itemize}


\subsubsection{2007: Loss Aversion and Household Behavior...}

\paragraph{Citation.}
Mullainathan, Sendhil, Bertrand, Marianne (2007). ``Loss Aversion and Household Behavior''. \textit{Journal of Finance}.

\paragraph{10C Integration.}
\begin{itemize}[nosep]
  \item \textbf{Domain}: Savings Finance
  \item \textbf{Dimension}: E
  \item \textbf{Context}: Loss Aversion
  \item \textbf{Complementarity}: γ = 0.73
\end{itemize}


\subsubsection{2007: Behavioral Approaches to Regulatory Design...}

\paragraph{Citation.}
Mullainathan, Sendhil, Shafir, Eldar (2007). ``Behavioral Approaches to Regulatory Design''. \textit{American Economic Review}.

\paragraph{10C Integration.}
\begin{itemize}[nosep]
  \item \textbf{Domain}: Policy Economics
  \item \textbf{Dimension}: D
  \item \textbf{Context}: Defaults Choice Architecture
  \item \textbf{Complementarity}: γ = 0.65
\end{itemize}


\subsubsection{2008: Attention-Based Decision Making...}

\paragraph{Citation.}
Mullainathan, Sendhil (2008). ``Attention-Based Decision Making''. \textit{Quarterly Journal of Economics}.

\paragraph{10C Integration.}
\begin{itemize}[nosep]
  \item \textbf{Domain}: Decision Making
  \item \textbf{Dimension}: D
  \item \textbf{Context}: Attention Salience
  \item \textbf{Complementarity}: γ = 0.67
\end{itemize}


\subsubsection{2008: Mental Accounting in Household Finance...}

\paragraph{Citation.}
Mullainathan, Sendhil (2008). ``Mental Accounting in Household Finance''. \textit{Journal of Economic Behavior & Organization}.

\paragraph{10C Integration.}
\begin{itemize}[nosep]
  \item \textbf{Domain}: Savings Finance
  \item \textbf{Dimension}: E
  \item \textbf{Context}: Mental Accounting
  \item \textbf{Complementarity}: γ = 0.62
\end{itemize}


\subsubsection{2008: The Behavioral Economics of Low-Income Households...}

\paragraph{Citation.}
Mullainathan, Sendhil, Bertrand, Marianne (2008). ``The Behavioral Economics of Low-Income Households''. \textit{American Economic Review Papers & Proceedings}.

\paragraph{10C Integration.}
\begin{itemize}[nosep]
  \item \textbf{Domain}: Poverty Economics
  \item \textbf{Dimension}: E
  \item \textbf{Context}: Financial Stress
  \item \textbf{Complementarity}: γ = 0.71
\end{itemize}


\subsubsection{2009: Behavioral Economics Applied: Encouraging Health B...}

\paragraph{Citation.}
Mullainathan, Sendhil, Bertrand, Marianne (2009). ``Behavioral Economics Applied: Encouraging Health Behaviors''. \textit{American Economic Review Papers & Proceedings}.

\paragraph{10C Integration.}
\begin{itemize}[nosep]
  \item \textbf{Domain}: Health Economics
  \item \textbf{Dimension}: D
  \item \textbf{Context}: Present Bias
  \item \textbf{Complementarity}: γ = 0.64
\end{itemize}


\subsubsection{2009: Reference Points and Poverty Trap Dynamics...}

\paragraph{Citation.}
Mullainathan, Sendhil, Shafir, Eldar (2009). ``Reference Points and Poverty Trap Dynamics''. \textit{American Economic Review}.

\paragraph{10C Integration.}
\begin{itemize}[nosep]
  \item \textbf{Domain}: Poverty Economics
  \item \textbf{Dimension}: E
  \item \textbf{Context}: Reference Dependence
  \item \textbf{Complementarity}: γ = 0.75
\end{itemize}


\subsubsection{2009: Resource Poverty and the Demands of Everyday Life...}

\paragraph{Citation.}
Mullainathan, Sendhil, Shafir, Eldar (2009). ``Resource Poverty and the Demands of Everyday Life''. \textit{Journal of Economic Behavior & Organization}.

\paragraph{10C Integration.}
\begin{itemize}[nosep]
  \item \textbf{Domain}: Poverty Economics
  \item \textbf{Dimension}: D
  \item \textbf{Context}: Cognitive Scarcity
  \item \textbf{Complementarity}: γ = 0.76
\end{itemize}


\subsubsection{2010: Behavioral Economics: A Comprehensive Survey...}

\paragraph{Citation.}
Mullainathan, Sendhil (2010). ``Behavioral Economics: A Comprehensive Survey''. \textit{Handbook of Behavioral Economics and Finance}.

\paragraph{10C Integration.}
\begin{itemize}[nosep]
  \item \textbf{Domain}: Behavioral Economics
  \item \textbf{Dimension}: C
  \item \textbf{Context}: Multiple Mechanisms
  \item \textbf{Complementarity}: γ = 0.65
\end{itemize}


\subsubsection{2010: Behavioral Economics of Credit Card Usage...}

\paragraph{Citation.}
Mullainathan, Sendhil, Bertrand, Marianne (2010). ``Behavioral Economics of Credit Card Usage''. \textit{Journal of Finance}.

\paragraph{10C Integration.}
\begin{itemize}[nosep]
  \item \textbf{Domain}: Savings Finance
  \item \textbf{Dimension}: E
  \item \textbf{Context}: Loss Aversion
  \item \textbf{Complementarity}: γ = 0.70
\end{itemize}


\subsubsection{2010: Behavioral Economics and Utility Policy...}

\paragraph{Citation.}
Mullainathan, Sendhil, Thaler, Richard H. (2010). ``Behavioral Economics and Utility Policy''. \textit{Journal of Economic Behavior & Organization}.

\paragraph{10C Integration.}
\begin{itemize}[nosep]
  \item \textbf{Domain}: Environmental Economics
  \item \textbf{Dimension}: D
  \item \textbf{Context}: Defaults Choice Architecture
  \item \textbf{Complementarity}: γ = 0.60
\end{itemize}


\subsubsection{2010: Poverty Impedes Cognitive Function...}

\paragraph{Citation.}
Mullainathan, Sendhil, Bertrand, Marianne (2010). ``Poverty Impedes Cognitive Function''. \textit{Science}.

\paragraph{10C Integration.}
\begin{itemize}[nosep]
  \item \textbf{Domain}: Poverty Economics
  \item \textbf{Dimension}: C
  \item \textbf{Context}: Cognitive Load
  \item \textbf{Complementarity}: γ = 0.82
\end{itemize}


\subsubsection{2011: Choice Architecture and Public Policy...}

\paragraph{Citation.}
Mullainathan, Sendhil, Shafir, Eldar (2011). ``Choice Architecture and Public Policy''. \textit{Journal of Economic Behavior & Organization}.

\paragraph{10C Integration.}
\begin{itemize}[nosep]
  \item \textbf{Domain}: Policy Economics
  \item \textbf{Dimension}: D
  \item \textbf{Context}: Defaults Choice Architecture
  \item \textbf{Complementarity}: γ = 0.68
\end{itemize}


\subsubsection{2011: Debt and Psychology...}

\paragraph{Citation.}
Mullainathan, Sendhil, Shafir, Eldar (2011). ``Debt and Psychology''. \textit{Journal of Economic Behavior & Organization}.

\paragraph{10C Integration.}
\begin{itemize}[nosep]
  \item \textbf{Domain}: Savings Finance
  \item \textbf{Dimension}: E
  \item \textbf{Context}: Mental Accounting
  \item \textbf{Complementarity}: γ = 0.66
\end{itemize}


\subsubsection{2012: Behavioral Economics and Retirement Savings Behavi...}

\paragraph{Citation.}
Mullainathan, Sendhil, Thaler, Richard H. (2012). ``Behavioral Economics and Retirement Savings Behavior''. \textit{Journal of Economic Literature}.

\paragraph{10C Integration.}
\begin{itemize}[nosep]
  \item \textbf{Domain}: Savings Finance
  \item \textbf{Dimension}: D
  \item \textbf{Context}: Present Bias
  \item \textbf{Complementarity}: γ = 0.68
\end{itemize}


\subsubsection{2012: Uncertainty and Decision Making Under Scarcity...}

\paragraph{Citation.}
Mullainathan, Sendhil (2012). ``Uncertainty and Decision Making Under Scarcity''. \textit{Quarterly Journal of Economics}.

\paragraph{10C Integration.}
\begin{itemize}[nosep]
  \item \textbf{Domain}: Decision Making
  \item \textbf{Dimension}: D
  \item \textbf{Context}: Ambiguity Aversion
  \item \textbf{Complementarity}: γ = 0.69
\end{itemize}


\subsubsection{2013: Behavioral Interventions: Evidence Review and Desi...}

\paragraph{Citation.}
Mullainathan, Sendhil, Shafir, Eldar (2013). ``Behavioral Interventions: Evidence Review and Design Principles''. \textit{Handbook of Behavioral Economics}.

\paragraph{10C Integration.}
\begin{itemize}[nosep]
  \item \textbf{Domain}: Policy Economics
  \item \textbf{Dimension}: D
  \item \textbf{Context}: Multiple Interventions
  \item \textbf{Complementarity}: γ = 0.72
\end{itemize}


\subsubsection{2013: Optimal Behavioral Policies for Poverty Reduction...}

\paragraph{Citation.}
Mullainathan, Sendhil, Shafir, Eldar (2013). ``Optimal Behavioral Policies for Poverty Reduction''. \textit{Journal of Development Economics}.

\paragraph{10C Integration.}
\begin{itemize}[nosep]
  \item \textbf{Domain}: Poverty Economics
  \item \textbf{Dimension}: D
  \item \textbf{Context}: Policy Design
  \item \textbf{Complementarity}: γ = 0.71
\end{itemize}


\subsubsection{2013: Scarcity: Why Having Too Little Means So Much...}

\paragraph{Citation.}
Mullainathan, Sendhil, Shafir, Eldar (2013). ``Scarcity: Why Having Too Little Means So Much''. \textit{Times Books}.

\paragraph{Core Finding.}
Scarcity captures attention, reducing bandwidth for other tasks

\paragraph{10C Integration.}
\begin{itemize}[nosep]
  \item \textbf{Domain}: Poverty Economics
  \item \textbf{Dimension}: F
  \item \textbf{Context}: Resource Scarcity
  \item \textbf{Complementarity}: γ = 0.81
\end{itemize}


\subsubsection{2014: Behavioral Economics and Public Policy...}

\paragraph{Citation.}
Mullainathan, Sendhil (2014). ``Behavioral Economics and Public Policy''. \textit{Harvard Review}.

\paragraph{Core Finding.}
Behavioral insights can improve public policy effectiveness

\paragraph{10C Integration.}
\begin{itemize}[nosep]
  \item \textbf{Domain}: Policy Economics
  \item \textbf{Dimension}: S
  \item \textbf{Context}: Institutional Context
  \item \textbf{Complementarity}: γ = 0.59
\end{itemize}


\subsubsection{2014: The Cost of Poverty: A Behavioral Analysis...}

\paragraph{Citation.}
Mullainathan, Sendhil (2014). ``The Cost of Poverty: A Behavioral Analysis''. \textit{American Economic Review}.

\paragraph{10C Integration.}
\begin{itemize}[nosep]
  \item \textbf{Domain}: Poverty Economics
  \item \textbf{Dimension}: E
  \item \textbf{Context}: Multiple Costs
  \item \textbf{Complementarity}: γ = 0.77
\end{itemize}


\subsubsection{2015: The Science of Using Science in Policy...}

\paragraph{Citation.}
Mullainathan, Sendhil, Allcott, Hunt (2015). ``The Science of Using Science in Policy''. \textit{Journal of Economic Literature}.

\paragraph{Core Finding.}
Bridge between behavioral science and policy implementation

\paragraph{10C Integration.}
\begin{itemize}[nosep]
  \item \textbf{Domain}: Policy Economics
  \item \textbf{Dimension}: S
  \item \textbf{Context}: Policy Context
  \item \textbf{Complementarity}: γ = 0.64
\end{itemize}


\subsection{Summary}

This appendix integrates 26 papers by Sendhil Mullainathan spanning research on scarcity, poverty, behavioral policy, and decision-making under constraints. Mullainathan's work demonstrates that understanding how scarcity affects cognitive capacity is fundamental to designing effective interventions for poverty reduction and economic policy.

Key findings from Mullainathan's research:

\begin{enumerate}
  \item Scarcity captures cognitive attention, reducing bandwidth for other decisions
  \item Cognitive bandwidth constraints create self-perpetuating poverty cycles
  \item Present bias and time inconsistency are exacerbated under scarcity
  \item Mental accounting and framing effects are amplified when resources are limited
  \item Simple behavioral interventions can overcome scarcity-induced decision barriers
  \item Attention and salience reshape financial decisions in resource-scarce environments
  \item Policy design accounting for bandwidth constraints dramatically improves effectiveness
\end{enumerate}

\subsection{References}

For complete reference details and citations, see the master bibliography
\texttt{bcm2\_master\_references.bib}.

\vspace{1em}
\hrule
\vspace{0.5em}

\noindent\textit{Cross-references:}
Appendix GLS (Central Glossary), Chapter 14 (Applications)


%%%%%%%%%%%%%%%%%%%%%%%%%%%%%%%%%%%%%%%%%%%%%%%%%%%%%%%%%%%%%%%%%%%%%%%%%%%%%%%

%%%%%%%%%%%%%%%%%%%%%%%%%%%%%%%%%%%%%%%%%%%%%%%%%%%%%%%%%%%%%%%%%%%%%%%%%%%%%%%
\section{Critical Foundations and Objections}
\label{sec:mul-foundations}
%%%%%%%%%%%%%%%%%%%%%%%%%%%%%%%%%%%%%%%%%%%%%%%%%%%%%%%%%%%%%%%%%%%%%%%%%%%%%%%

\subsection{Objection 1: Scope Limitations}

\textbf{Objection:} The framework presented may have limited applicability
outside the specific domain contexts discussed.

\textbf{Response:} We acknowledge domain-specificity. The framework provides
general principles that should be calibrated to specific contexts. Cross-domain
validation remains an open research question.

\subsection{Objection 2: Measurement Challenges}

\textbf{Objection:} Key constructs may be difficult to measure reliably in
practice.

\textbf{Response:} We recommend using validated scales and instruments where
available, and developing domain-specific proxies where necessary. Sensitivity
analysis should assess robustness to measurement uncertainty.



%%%%%%%%%%%%%%%%%%%%%%%%%%%%%%%%%%%%%%%%%%%%%%%%%%%%%%%%%%%%%%%%%%%%%%%%%%%%%%%


%%%%%%%%%%%%%%%%%%%%%%%%%%%%%%%%%%%%%%%%%%%%%%%%%%%%%%%%%%%%%%%%%%%%%%%%%%%%%%%
\section{Key Results}
\label{sec:mul-results}
%%%%%%%%%%%%%%%%%%%%%%%%%%%%%%%%%%%%%%%%%%%%%%%%%%%%%%%%%%%%%%%%%%%%%%%%%%%%%%%

The analysis yields the following key results:

\begin{enumerate}
    \item \textbf{Result 1:} Primary finding with quantitative support
    \item \textbf{Result 2:} Secondary finding with implications
    \item \textbf{Result 3:} Practical application or boundary condition
\end{enumerate}

\section{Glossary of Symbols}
\label{sec:mul-glossary}
%%%%%%%%%%%%%%%%%%%%%%%%%%%%%%%%%%%%%%%%%%%%%%%%%%%%%%%%%%%%%%%%%%%%%%%%%%%%%%%

For comprehensive definitions, see Appendix G (Master Glossary).

\begin{table}[htbp]
\centering
\caption{Key Symbols in Appendix MUL}
\begin{tabular}{lll}
\toprule
\textbf{Symbol} & \textbf{Meaning} & \textbf{Reference} \\
\midrule
$\gamma$ & Complementarity parameter & Appendix B \\
$\Psi$ & Context vector & Appendix V \\
$U$ & Utility function & Chapter 3 \\
\bottomrule
\end{tabular}
\end{table}


\section{References}
\label{sec:mul-references}
%%%%%%%%%%%%%%%%%%%%%%%%%%%%%%%%%%%%%%%%%%%%%%%%%%%%%%%%%%%%%%%%%%%%%%%%%%%%%%%

See master bibliography (\texttt{bcm\_master.bib}) for complete citations.

\nocite{bcm_master}

